
This section details experimental protocols for the three sub-hypotheses, including datasets, evaluation metrics, and success criteria. The design follows the decomposition strategy (Section 3): test data collection (h-e1), then correlation (h-m1), then training (h-m2) sequentially, with each stage building on the prior.

\textbf{H-E1: Benchmark Data Collection Experiment.} The research question is whether at least 50 benchmarks with complete method comparisons can be collected from published sources. Data sources include OGB (Open Graph Benchmark with 15 graph datasets and documented baselines for GCN, GraphSAINT, ClusterGCN, GraphSAGE families, accessible via Python library pip install ogb), FedML (6 federated learning datasets with FedAvg, FedProx, FedOpt baselines, accessible via GitHub repository), LEAF (5 datasets with heterogeneity statistics and baseline results, available via GitHub), pFL-Bench (8 personalized federated learning benchmarks, accessible via Papers with Code or direct repository links), Champneys et al. NLSI (5 nonlinear system identification tasks with NARX, LSTM, Polynomial baseline comparisons, results reported in paper tables), and Zhou et al. Medical FL (9 medical imaging datasets with FedAvg, SCAFFOLD, FedDyn, augmentation method results, reported in paper).

The collection protocol follows: (1) API/library access---load datasets via official libraries (OGB, FedML) to verify accessibility and extract metadata from loaded objects, (2) GitHub repository---clone or fetch README files documenting dataset statistics and baseline results, (3) paper table extraction---for Champneys and Zhou, manually extract benchmark names, dataset characteristics if reported, and method rankings from paper tables into structured CSV format. Metadata fields per benchmark include Tier 1 features (universal, under one second: sample\_size n, dimensionality d, num\_classes for classification, class\_imbalance as Gini coefficient of class distribution), Tier 2 features (domain-specific, under one minute: autocorrelation\_lag1 for time-series, edge\_density for graphs, feature\_correlation\_rank for tabular), and method rankings (percentile ranking 0-100\% for representative methods from each family: Linear, Polynomial, RNN/LSTM, Augmentation).

Success criteria: PASS requires at least 50 benchmarks collected with at least 80\% of Tier 1 features populated (sample\_size, dimensionality, class\_imbalance available for at least 40 benchmarks). PARTIAL indicates 30-49 benchmarks collected or below 80\% feature completeness (data availability challenges). FAIL is fewer than 30 benchmarks collected (insufficient training data for meta-classifier; hypothesis untestable). Collected benchmarks should ensure generalization by spanning domains (vision image classification, NLP/text, tabular, time-series, graph), scales (small n<1K, medium $1K\leq n$<10K, large $n\geq 10K)$, task types (classification, regression), and structural properties (balanced vs. imbalanced classes, low-dimensional vs. high-dimensional features).

\textbf{H-M1: Feature-Method Correlation Analysis Experiment.} The research question is whether dataset characteristics correlate with method family performance, validating that fast-to-compute features capture information relevant for method selection. Features are computed in two tiers: Tier 1 (universal, under one second) includes sample\_size (total training samples), dimensionality (number of input features), num\_classes (number of output classes for classification), class\_imbalance (Gini coefficient of class distribution, 0=balanced, 1=single class dominates), and signal\_noise\_ratio (variance of features divided by variance of residuals for regression). Tier 2 (domain-specific, 5-15 seconds) includes autocorrelation\_lag1 (first-order autocorrelation for time-series), edge\_density (average degree divided by max degree for graphs), and feature\_correlation\_rank (rank of feature correlation matrix for tabular). Tier 1 features should be available for all benchmarks (universal), while Tier 2 features are computed conditionally based on domain.

Method family representatives map diverse methods to four broad families based on inductive bias: Linear (ARX, Ridge Regression, Linear SVM), Polynomial (NARX-Poly, Polynomial Regression, Kernel SVM with polynomial kernel), RNN (LSTM, GRU, vanilla RNN), and Augmentation (data augmentation techniques: DDPM+LS, MixUp, CutMix). For each benchmark, the ranking percentile (0-100\%) of the best-performing method in each family is extracted. If multiple representatives exist (e.g., LSTM and GRU), the maximum ranking within the family is used.

The correlation protocol, for each (feature, method\_family) pair, involves: (1) data preparation---filter benchmarks to those with non-missing values for both feature and method ranking, (2) correlation test---compute Spearman's rank correlation $\rho$ (robust to outliers and nonlinear monotonic relationships), (3) significance test---report p-value with threshold p<0.05 for significance, (4) practical significance---require |$\rho$|>0.3 for meaningful correlation (weak correlations uninformative for prediction). All combinations are tested: 7 features (5 Tier 1 + 2 domain-conditional Tier 2) $\times$ 4 method families = 28 tests. Bonferroni correction would require p<0.05/$28\approx 0.0018$, but uncorrected p-values are reported (exploratory analysis, not confirmatory hypothesis testing).

Expected patterns from literature include: sample\_size $\leftrightarrow$ Augmentation (negative correlation: small datasets benefit more from augmentation), dimensionality $\leftrightarrow$ Linear (positive: high-dimensional favors regularized linear models), autocorrelation\_lag1 $\leftrightarrow$ RNN (positive: temporal structure favors recurrent models), and class\_imbalance $\leftrightarrow$ Augmentation (positive: imbalanced datasets benefit from synthetic samples).

Success criteria: PASS requires at least 5 significant correlations ($\rho$>0.3, p<0.05) spanning at least 3 method families (evidence that features capture method-relevant information). PARTIAL indicates 3-4 significant correlations or correlations confined to single family (suggests feature insufficiency). FAIL is fewer than 3 significant correlations (features non-informative; either need Tier 3 expensive probing or dataset characteristics do not predict method performance).

\textbf{H-M2: Meta-Classifier Training and Evaluation Experiment.} The research question is whether a Random Forest trained on collected benchmarks can predict method family rankings on held-out datasets with at least 30\% top-30\% success rate. The model architecture uses Random Forest (scikit-learn RandomForestClassifier) for interpretability and robustness. Input is a 7-dimensional feature vector (Tier 1 + Tier 2 features, domain-conditional). Output is a 4-class prediction (Linear, Polynomial, RNN, Augmentation). Hyperparameters include n\_estimators=100 (ensemble size, standard default), max\_depth=10 (prevent overfitting on small training set), and min\_samples\_split=5 (require at least 5 samples to split node). Random Forest is interpretable via SHAP feature importance, handles missing values (NaN indicators), and is robust to small sample sizes.

The training protocol defines the target variable per benchmark as the best-performing method family (highest ranking percentile among the four families), converting regression (predicting exact ranking) to 4-class classification (predicting family). Cross-validation uses leave-5-out CV with 10 folds (if 50 benchmarks collected, each fold trains on 45 and tests on 5), stratified by domain if possible (ensure each fold has vision/NLP/tabular/time-series representation). Evaluation metrics include primary (top-30\% success rate---percentage of held-out benchmarks where predicted family achieves at most 30th percentile ranking) and secondary (confusion matrix, per-domain accuracy, feature importance via SHAP).

Baseline comparisons include: random selection (uniform random choice from Linear, Polynomial, RNN, Augmentation, expected 30\% top-30\% success rate by definition), majority class (always predict most frequent winner in training set, expected 30-40\% success rate as degeneracy check), and domain folklore (predict based on domain only: $vision\rightarrow RNN$ under pre-Transformer era assumption, time-$series\rightarrow RNN, tabular\rightarrow Linear$, expected 40-50\% success rate per Phase 2A estimates).

Success criteria: PASS requires at least 35\% top-30\% success rate, significantly better than random (Chi-square test p<0.05). PARTIAL indicates 30-35\% success rate, marginally better than random (suggests weak but present signal; investigate feature sufficiency or increase training data). FAIL is below 30\% success rate (no better than random) or worse than majority-class baseline (indicates no learning; either features insufficient, sample size too small, or method rankings too chaotic).
